\documentclass[11pt]{article}
\usepackage[utf8]{inputenc}
\usepackage[T1]{fontenc}
\usepackage{amsmath}
\usepackage{amsfonts}
\usepackage{amssymb}
\usepackage{geometry}
\usepackage{listings}
\usepackage{xcolor}

\geometry{letterpaper, margin=1in}

\definecolor{codegreen}{rgb}{0,0.6,0}
\definecolor{codegray}{rgb}{0.5,0.5,0.5}
\definecolor{codepurple}{rgb}{0.58,0,0.82}
\definecolor{backcolour}{rgb}{0.95,0.95,0.92}

\lstdefinestyle{mystyle}{
    backgroundcolor=\color{backcolour},   
    commentstyle=\color{codegreen},
    keywordstyle=\color{magenta},
    numberstyle=\tiny\color{codegray},
    stringstyle=\color{codepurple},
    basicstyle=\ttfamily\footnotesize,
    breakatwhitespace=false,         
    breaklines=true,                 
    captionpos=b,                    
    keepspaces=true,                 
    numbers=left,                    
    numbersep=5pt,                  
    showspaces=false,                
    showstringspaces=false,
    showtabs=false,                  
    tabsize=2
}

\lstset{style=mystyle}

\title{\textbf{Editorial: Problem F - Optimal Mirroring}}
\author{Setter: Riverxia}
\date{}

\begin{document}

\maketitle

\section*{Problem Overview}
Touko wants to maintain the same relative order as a fixed permutation $a$ for $q$ different intervals. Yuu needs to find the minimum number of edges in a DAG $G_k$ such that its satisfying permutations are exactly those that match Touko's sister's relative order on the first $k$ intervals. This is equivalent to finding the size of the \textbf{transitive reduction} of the partial order defined by the union of constraints from the $k$ intervals.

\section*{Solution Sketch}

\subsection*{Key Observations}
\begin{enumerate}
    \item An edge $(u, v)$ is required in the transitive reduction if and only if $u < v$ in the partial order and there is no path of length $\ge 2$ from $u$ to $v$.
    \item For a single interval $[l, r]$, the minimal set of edges is simply $x \to y$ for all pairs $(x, y)$ such that $a_x$ and $a_y$ are \textbf{consecutive values} in the sorted sub-permutation $a_l, \dots, a_r$. There are exactly $(r-l)$ such edges.
    \item When multiple intervals overlap, an edge $u \to v$ from one interval might be rendered redundant by a path in another.
\end{enumerate}

\subsection*{The Algorithm}
A crucial insight from the original problem (CF 1508F) is that we only need to consider ``candidate edges.'' These are edges $(u, v)$ such that $u$ and $v$ are adjacent in value in \textit{some} range.

Using Mo's algorithm, we can maintain the set of edges where the values are adjacent. As we move the pointers $l$ and $r$, we add and remove edges. Each time an edge $(u, v)$ is formed, we track the range of $k$ for which it exists.

However, many of these edges are redundant. Lemma 2 from the tutorial states that for any vertex $u$, there is at most one ``right edge'' $u \to u_r$ and at most one ``left edge'' $u \to u_l$ in the transitive reduction at any time.

We can identify these candidate edges by running Mo's algorithm. There are $O(N \sqrt{Q})$ such candidate edges. For each candidate edge, we determine the time interval $[t_{start}, t_{end}]$ during which it is the ``best'' candidate and not ``hidden'' by a path of length 2.

\section*{Complexity Analysis}
\begin{itemize}
    \item \textbf{Time Complexity:} $O(N \sqrt{Q} \log N)$ or $O(N \sqrt{Q})$ depending on the implementation of the candidate edge tracking. Given the 7s time limit, a Mo's algorithm approach is feasible.
    \item \textbf{Space Complexity:} $O(N + Q)$ to store the permutation, queries, and candidate edge info.
\end{itemize}

\section*{Bloom into You Connection}
Touko's struggle is about the pressure of maintaining a specific relative order (her sister's legacy). Yuu's task of finding the \textit{minimum} number of edges represents her finding the most fundamental truths about Touko—stripping away the redundant layers of the performance until only the essential constraints remain.

\end{document}